\documentclass[10pt,a4paper]{amsart}
\usepackage[margin=19mm]{geometry}
\usepackage{amsmath,amssymb,mathtools}
\usepackage[scr=rsfs,cal=euler]{mathalfa}
\usepackage{xfrac}
\usepackage[unicode,hidelinks,bookmarks=false]{hyperref}
\newcommand{\ie}{i.e.\ }
\newcommand{\cf}{cf.\ }
\newcommand{\resp}{respectively\ }
\newcommand{\Subsection}{\subsection}
\providecommand{\nobreakdash}{\nobreak\hbox{-}}
\setlength{\emergencystretch}{2em}
\multlinegap0pt
\usepackage{bm}
\usepackage{bbm}
\usepackage{dsfont}
\usepackage{xspace}
\usepackage{cancel}
\usepackage{mathtools}
\usepackage{cleveref}
\usepackage{amsthm}
\makeatletter
\@ifundefined{theor}{\newtheorem{theor}{Theor}}{}
\@ifundefined{lemma}{\newtheorem{lemma}[theor]{Lemma}}{}
\@ifundefined{prop}{\newtheorem{prop}[theor]{Proposition}}{}
\@ifundefined{rem}{\newtheorem{rem}[theor]{Remark}}{}
\makeatother
\DeclareMathOperator{\End}{End}
\title[Set-theoretic solutions of the Yang-Baxter equation from inv]{Set-theoretic solutions of the Yang-Baxter equation from inverse braces}
\author{Francesco Catino, Marzia Mazzotta, Susanna Tieni}
\date{Source: arXiv:2607.16045v1 (CC BY 4.0)}
\begin{document}
\maketitle

\noindent\emph{Source and licence.} Set-theoretic solutions of the Yang-Baxter equation from inverse braces. Version arXiv:2607.16045v1 (CC BY 4.0), distributed under the Creative Commons Attribution 4.0 International licence, as recorded in the arXiv licence metadata for this exact version and shown on the abstract page https://arxiv.org/abs/2607.16045v1. Full rights notice: licenses/arxiv-2607.16045-CC-BY-4.0.txt.\par\smallskip

\noindent\textit{Editorial scope.} Counted unit: the Theor whose source label is \texttt{teo\_solu} in the article (source0.tex lines 951--959, source label \texttt{teo\_solu}), delivered together with the article's own complete original proof of it (source0.tex lines 960--1009, one \texttt{proof} environment of 50 lines). No mathematics is written, repaired or re-derived: statement and proof are the article's own text line for line. Required context retained uncounted: key2.4 (lines 462--464); lemma\_ab (lines 468--474); lemma\_1 (lines 557--563); rem\_sigma\_lambda (lines 665--686); A (lines 780--782); A* (lines 818--827); sigma\_rho (lines 859--861); condizione\_a\_2 (lines 897--908); B (lines 898--901). Every definition, display and label of those lines is retained unchanged. Omitted from this package and not redistributed: 1--461, 465--467, 475--556, 564--664, 687--779, 783--817, 828--858, 862--896, 902--950, 1010--1810. Nothing inside a retained excerpt is omitted or altered: every retained body is delivered exactly as the article's lines are written, so no omission receipt of any kind accompanies this package. Citations: the retained excerpts carry no citation command at all, so no bibliography entry of the article is transcribed and no reference is redistributed. Reference adaptation: none was needed and none was made; every reference inside the delivered unit resolves inside it, so no unresolved reference or citation is produced. Printed slips kept literal: no printed word, symbol, hypothesis or number of the counted statement or of its proof was altered. Layout and class: the article's own class is \texttt{elsarticle} with packages including lipsum, hyperref, enumitem, circuitikz, tikz-cd, multirow, url, pst-node, comment, cancel, soul, tikz, latexsym, indentfirst; the wrapper is the lane's pinned \texttt{amsart} editorial wrapper at 10pt on A4, which restates the theorem-like environments the retained text needs and adds the article's own macro definitions that the retained text needs (\texttt{\textbackslash End}), with extra wrapper packages (bm, bbm, dsfont, xspace, cancel, mathtools, cleveref). The delivered numbering is the unit's own. Assets: no retained span contains a figure environment, a graphics-inclusion command, a PDF-page inclusion, a table or a TikZ picture; no image file is copied into this package, the package declares no asset dependency and no image file is redistributed with it. Licence: version arXiv:2607.16045v1 of this article is distributed under the Creative Commons Attribution 4.0 International licence, as recorded in the arXiv licence metadata for this exact version and shown on the abstract page https://arxiv.org/abs/2607.16045v1. A full rights notice is delivered at licenses/arxiv-2607.16045-CC-BY-4.0.txt. Characters: every retained excerpt is pure ASCII, so the delivered engine, XeLaTeX with its default Latin Modern text font, reports no missing character.
\begin{lemma}\label{key2.4}
Let $(S,+,\circ)$ be an inverse brace. Then, for every $a\in S$, $\lambda_a \in \End(S,+)$. Moreover, the map $\lambda: S \rightarrow \End(S,+),\; a \mapsto \lambda_a$ is a homomorphism from the semigroup $(S,\circ)$ to $\End(S,+)$.
\end{lemma}

\begin{lemma}\label{lemma_ab}
Let $(S,+,\circ)$ be an inverse brace. Then, for all $a,b\in S$, the following identity hold:
\begin{enumerate}
\item[\rm{(1)}] $a\circ b= a - a\circ (-b)+a$,
\item[\rm{(2)}] $a\circ b = a+\lambda_a(b)$.
\end{enumerate}
\end{lemma}

\begin{lemma}\label{lemma_1}
    Let $(S,+,\circ)$ be an inverse brace. Then, for all $a,b\in S$, the following identity hold: 
\begin{enumerate}
\item[\rm{(1)}] $a\circ b = a\circ b -a+a,$
\item[\rm{(2)}] $a= a -a \circ a^-+a\circ a^-$ .
\end{enumerate}
\end{lemma}  

\begin{prop}\label{rem_sigma_lambda}
If $(S, +, \circ)$ is an inverse brace, then, for all $a, b \in S$,
$$\sigma_a(b)=a \circ a^- + \lambda_a(b)=a \circ a^- \circ \lambda_a(b).$$
Moreover, $\sigma_a=\lambda_a$, for every $a \in S$, if and only if $S$ is a weak brace.
    \begin{proof}
        Let $a, b \in S$. Then, by \cref{key2.4}, we have
        \begin{align*}
        \sigma_a(b)&=a \circ a^-+\lambda_a( b)=a \circ a^-+\lambda_{a \circ a^- \circ a}(b)=a \circ a^-+\lambda_{a \circ a^- }\lambda_a(b)=a \circ a^- \circ \lambda_a(b),
        \end{align*}
        where we apply \cref{lemma_ab}(2) in the last equality. \\
   Now, assume that $\sigma_a=\lambda_a$, for every $a \in S$. If $a \in S$, %Observing that $\lambda_a\left(a^- \circ a\right)=-a+a$, 
 we have
   \begin{align*}
      -a+a&=\lambda_a\left(a^- \circ a\right)\\
      &=\sigma_a\left(a^- \circ a\right)\\
      &=a \circ a^- + \lambda_a\left(a^- \circ a\right)\\
      &=a \circ a^--a+a\\
      &=a \circ a^- &\mbox{by \cref{lemma_1}(1)}
   \end{align*}
   Hence, $S$ is a weak brace. The converse implication is obvious.
    \end{proof}
\end{prop}

   \begin{align}\label{A} 
    &\sigma_a(e)=a \circ e \circ a^-\tag{I}
     \end{align}

\begin{rem}\label{A*}
    If $(S, +, \circ)$ is an inverse brace satisfying \eqref{A}, we also have that $\sigma_a(b)=a \circ b \circ b^- \circ \left(a^-+b\right)$, for all $a,b \in S$. In fact, 
   \begin{align*}
       a \circ b \circ b^- \circ \left(a^-+b\right)&=a \circ b \circ b^- \circ a^--a \circ b \circ b^-+a \circ b\\
       &=\sigma_a\left(b \circ b^-\right)-a \circ b \circ b^-+a \circ b &\mbox{by \eqref{A}}\\
       &=a \circ a^--a+a \circ b \circ b^--a \circ b \circ b^-+a \circ b &\mbox{by \cref{rem_sigma_lambda}}\\
       &=a \circ a^--a+a \circ b &\mbox{by \cref{lemma_ab}(2)}\\
       &=\sigma_a(b).
   \end{align*}
\end{rem}

\begin{lemma}\label{sigma_rho}
Let $(S,+,\circ)$ be an inverse brace such that \eqref{A} is satisfied. Then, for all $a, b \in S$, we have $\sigma_a(b)=a \circ b \circ \rho_b(a)^-.$
\end{lemma}

\begin{prop}\label{condizione_a_2}
  Let $(S, +, \circ)$ be an inverse brace satisfying \eqref{A} and such that \begin{align}
    &\sigma_a(b)\circ \sigma_a(b)^-=\sigma_a\left(b \circ b^-\right) 
    \tag{II}\label{B}
    \end{align}
    holds, for all $a, b \in S$. Then, the following hold:
    \begin{enumerate}
        \item[\rm{(1)}] $\sigma_a(b)\circ \rho_b(a)=a \circ b$, for all $a, b \in S$;
        \item[\rm{(2)}] the map $\rho: S \to S^S, b \mapsto \rho_b$
    is an anti-homomorphism from the inverse semigroup $(S, \circ)$ into the monoid $S^S$.
    \end{enumerate}
\end{prop}

\begin{theor}\label{teo_solu}
  Let $(S, +, \circ)$ be an inverse brace such that the map $\sigma: S \to S^S, a \mapsto \sigma_a$ is a homomorphism from the inverse semigroup $(S, \circ)$ to the monoid $S^S$, namely \eqref{A} is satisfied. Then the map $r_S:S \times S\to S\times S$ defined by $$r_S(a,b)=\left(\sigma_a(b), \sigma_a(b)^- \circ a \circ b\right),$$
 for all $a, b \in S$, is a solution if and only if the condition
    \begin{align}
    &\sigma_a(b)\circ \sigma_a(b)^-=\sigma_a\left(b \circ b^-\right), 
    \tag{II}
    \end{align}
    is satisfied, for all $a, b \in S$.
\end{theor}

\begin{proof}
First, assume that \eqref{B} holds. It is a routine computation to check that the map $r_S(a,b)=(\sigma_a(b), \rho_b(a))$ is a solution if and only if the following three fundamental equalities hold
\begin{align}
     &\label{first} \sigma_a\sigma_b(c)=\sigma_{\sigma_a\left(b\right)}\sigma_{\rho_b\left(a\right)}\left(c\right)\tag{Y1},\\
    &  \label{second}\sigma_{\rho_{\sigma_b\left(c\right)}\left(a\right)}\rho_c\left(b\right)=\rho_{\sigma_{\rho_b\left(a\right)}\left(c\right)}\sigma_a\left(b\right)\tag{Y2},\\
      &\label{third}\rho_c\rho_b(a)=\rho_{\rho_c\left(b\right)}\rho_{\sigma_b\left(c\right)}\left(a\right)\tag{Y3},
  \end{align} 
for all $a,b,c \in S$. \\
For $a, b, c \in S$, using \cref{condizione_a_2}, we obtain 
\begin{align*}
    \sigma_a\sigma_b(c)= \sigma_{a \circ b}(c)=\sigma_{\sigma_a(b)\circ \rho_b(a)}(c)=\sigma_{\sigma_a\left(b\right)}\sigma_{\rho_b\left(a\right)}\left(c\right)
\end{align*}
and
\begin{align*}
    \rho_c\rho_b(a)= \rho_{b \circ c}(a)=\rho_{\sigma_b(c)\circ \rho_c(b)}(a)=\rho_{\rho_c\left(b\right)}\rho_{\sigma_b\left(c\right)}\left(a\right).
\end{align*}
Hence, both \eqref{first} and \eqref{third} hold.
Moreover, 
\begin{align*}
    \rho_{\sigma_{\rho_b(a)}(c)}\sigma_a(b)&=\left(\sigma_{\sigma_a(b)}\sigma_{\rho_b(a)}(c)   \right)^-\circ \sigma_a(b)\circ\sigma_{\rho_b(a)}(c)
    %&\mbox{by \cref{condizione_a_2}(2)}
    \\&=\left(\sigma_{a}\sigma_b(c)\right)^- \circ \sigma_a(b)\circ \sigma_{\rho_b(a)}(c) &\mbox{by \eqref{first}}\\
    &=\left(\sigma_{a}\sigma_b(c)\right)^- \circ \sigma_a(b)\circ \rho_b(a)\circ c \circ\left(\rho_c\rho_b(a)\right)^-  &\mbox{by \cref{sigma_rho}}\\
    &=\left(\sigma_{a}\sigma_b(c)\right)^- \circ a \circ b \circ c \circ \left(\rho_c\rho_b(a)\right)^-  &\mbox{by \cref{condizione_a_2}(1)}\\
    %&=\left(\sigma_{a}\sigma_b(c)\right)^- \circ a \circ b \circ c \circ \left(\sigma_{\tau_b(a)}(c)^-\circ \tau_b(a) \circ c\right)^-  &\mbox{by \eqref{ab} \&\, \cref{condizione_a_2}-2.}\\
    %&=\left(\sigma_{a}\sigma_b(c)\right)^- \circ a\circ b \circ c \circ c^- \circ \tau_b(a)^- \circ \sigma_{\tau_b(a)}(c)\\
    &=\left(\sigma_{a}\sigma_b(c)\right)^- \circ a \circ \sigma_b(c) \circ \rho_c(b) \circ \left(\rho_c \rho_b(a)   \right)^-&\mbox{by \cref{condizione_a_2}(1)}\\
    &=\rho_{\sigma_b(c)}(a) \circ \rho_c(b) \circ \left(\rho_{\rho_c(b)}\rho_{\sigma_b(c)}(a)\right)^-&\mbox{by \eqref{third}}\\
&=\sigma_{\rho_{\sigma_b\left(c\right)}\left(a\right)}\rho_c\left(b\right). &\mbox{by \cref{sigma_rho}}
\end{align*}
Therefore, $r_S$ is a solution.\\
Conversely, assume that $r_S$ is a solution. From equation \eqref{first}, taking $c=-b^-+b^-$ we obtain
   \begin{align*}
\sigma_a\sigma_b\left(-b^-+b^-\right)=\sigma_{\sigma_a\left(b\right)}\sigma_{\rho_b\left(a\right)}\left(-b^-+b^-\right).  
   \end{align*}
   The left-hand side is equal to $
  \sigma_a\sigma_b\left(-b^-+b^-\right)=\sigma_a\left(b \circ \left( b^--b^-+b^-\right)\right)=\sigma_a\left(b \circ b^- \right)$. About the right-hand side, we have
  \begin{align*}
&\sigma_{\sigma_a\left(b\right)}\sigma_{\rho_b\left(a\right)}\left(-b^-+b^-\right)\\
   &=\sigma_{\sigma_a\left(b\right)}\left(\rho_b(a) \circ \left(b^- \circ a^- \circ \sigma_a(b)-b^-+b^-\right)\right)\\
   &=\sigma_{\sigma_a\left(b\right)}\left(\rho_b(a) \circ b^- \circ a^- \circ \sigma_a(b)\right) &\mbox{by \cref{lemma_1}(1)}\\
&=\sigma_{\sigma_a\left(b\right)}\left(\rho_b(a) \circ \rho_b(a)^-\right)\\
  &=\sigma_a(b) \circ \sigma_a(b)^--\sigma_a(b)+ \sigma_a(b) \circ \rho_b(a) \circ \rho_b(a)^-\\
    &=\sigma_a(b) \circ \sigma_a(b)^--\sigma_a(b)+ a \circ b \circ b^- \circ a^- \circ \sigma_a(b)\\
   &=\sigma_a(b) \circ \sigma_a(b)^--\sigma_a(b)+a \circ b \circ b^- \circ \left(a^-+b\right)\\
   &= \sigma_a(b) \circ \sigma_a(b)^--\sigma_a(b)+\sigma_a(b) &\mbox{by \cref{A*}}\\
   &=\sigma_a(b) \circ \sigma_a(b)^- &\mbox{by \cref{lemma_1}(1)}
  \end{align*}
Therefore, \eqref{B} is satisfied.
\end{proof}

\end{document}
